\section{Funktioner af flere variable}

\subsection{Kvadratisk Form}

Andengrads-ligninger af flere variable kan skrives som nedenstående, hvor $\vec x$ er en søjlevektor.

\[
    q: \mathbb{R}^n \to \mathbb{R}, \qquad q(\vec x) = \vec x^T \mtrix A \vec x + \vec x^T \vec b + c
    \kildetag{Definition 1.2.1}
\]

\bemærkBig{
    Transponering af kvadratisk form ($\vec x^T \mtrix A^T \vec x + \vec b^T \vec x + c$) ændrer ikke resultatet.
}

Matrix $\mtrix A$ er symmetrisk. 
\kilde{Lemma 2.8.5}

\subsection{Vektor-funktioner}
Vektor-funktioner tager et input $\mathbb R^n$ og giver et output i $\mathbb R^k$. For hver af $k$ indgange i resultatet kan en funktion skrives der betegner hvordan den indgang beregnes.
\[
    \vec f: \mathbb{R}^n \to \mathbb{R}^k, \qquad \vec f(\vec x) = \begin{bmatrix}
    f_1(\vec x) \\
    f_2(\vec x) \\
    \vdots \\
    f_k(\vec x)
    \end{bmatrix}
\]

\subsubsection{Vektor-felter}
En vektor-funktion er et vektor-felt $\iff \vec f: \dom(\vec f) \to \mathbb R^n, \quad \dom(\vec f) \subseteq \mathbb R^n$

Repræsenteres ved at tegne en vektorpil i punkter i domænet, hvor pilens retning og længde repræsenterer funktionens output i det punkt.

\subsection{Niveaukurver}
En niveaukurve for en en funktion og et niveau er mængden af alle værdier i domænet, hvor funktionen giver samme niveau.

\[
    \{\vec x \in \dom(\vec f) | \vec f(\vec x) = c\}
\]

\subsection{Nogle funktioner}
\subsubsection{ReLU}
ReLU er en aktiveringsfunktion. Aktiveringsfunktioner betegnes $\sigma$ og er ikke-lineære.
ReLU kan også bruges på vektorer, elementvist på hver indgang.
\[
\text{ReLU}(x) =
    \begin{cases}
    0 & \text{for } x < 0 \\
    x & \text{for } x \geq 0
    \end{cases}
\]

\subsubsection{SoftMax}
SoftMax er en aktiveringsfunktion, og en vektor-funktion. Den tager for hver ingang i vektoren $x_i$, og finder eksponentialet af den, og dividerer det med summen af eksponentialerne for alle indgange i vektoren:
\[
    \text{SoftMax}: \mathbb R^n \to \mathbb R^n, \qquad 
    \text{SoftMax}(x_i) = \frac{e^{x_i}}{\sum_{j=1}^n e^{x_j}} \quad
    \forall i \in \{1, 2, \ldots, n\}
\]
SoftMax giver en sandsynlighedsvektor.

\subsubsection{Netværksfunktion}
Netværksfunktionen $\Phi$ består af vægt-matricer $\mtrix A$, bias-vektorer $\vec b$ og aktiveringsfunktioner $\sigma$.
Funktionen laver altså først en affin afbildning ($\mtrix A x + \vec b$) og laver en ikke-lineær afbildning af resultatet.


\[
    \vec \Phi(\vec x) = \sigma(\mtrix A \cdot \vec x + \vec b)
\]
Netværksfunktioner kan laves med flere lag, f.eks.
\[
    \vec \Phi(\vec x) = \text{softMax}(\mtrix A_3 \text{ReLU}(\mtrix A_2 + \text{ReLU}(\mtrix A_1 \vec x + \vec b_1) \vec b_2) + \vec b_3)
\]